\documentclass[a4j,dvipdfmx]{jsarticle}
\usepackage{amsmath,amssymb}
\usepackage{siunitx}
\usepackage{ascmac}
\usepackage{amsthm}
\usepackage{bm}

% \usepackage[margin=10truemm]{geometry}

\DeclareMathOperator{\Span}{span}

\newcommand{\R}{\mathbb{R}}
\newcommand{\C}{\mathbb{C}}
\newcommand{\K}{\mathbb{K}}
\newcommand{\rank}{\mathrm{rank}}
\newcommand{\tr}{\mathrm{tr}}
\newcommand{\Ker}{\mathrm{Ker}\hspace{0.5mm}}
\newcommand{\E}{\mathcal{E}}
\newcommand{\F}{\mathcal{F}}
\newcommand{\trans}{{}^t}
\renewcommand{\labelenumi}{(\arabic{enumi})}
\renewcommand{\Im}{\mathrm{Im}\hspace{0.5mm}}
\newcommand{\Slash}{\hrulefill 演習問題中にあったもの \hrulefill}

\usepackage{enumitem}
\setlist[itemize]{before=\normalfont}

% 定理スタイルの設定
\newtheoremstyle{mystyle}% スタイル名
  {}% 上部スペース
  {}% 下部スペース
  {\normalfont}% 本文フォント（斜体。立体が良い場合は \normalfont）
  {}% インデント量
  {\bfseries}% 見出しフォント
  {.}% 見出し後の記号
  { }% 見出し後のスペース
  {\thmname{#1}\thmnumber{ #2}\thmnote{ (#3)}}% 見出しの書式

\theoremstyle{mystyle}

% 1. definition: 番号なし（通し番号をなくす）
\newtheorem*{definition}{定義}

% 2. theorem と proposition: 番号を共有し、セクションごとにリセット
% [section] をつけることで 5.1, 5.2 の形式になります
\newtheorem{theorem}{定理}[section]
\newtheorem{proposition}[theorem]{命題}

% 3. 性質: 番号なし
\newtheorem*{property}{性質}

\begin{document}
\part*{線型代数ゼミ 要点}
\begin{itembox}{記号について}
    Review 以外は「定義(Def)」「定理(Thm)」「命題(Prop)」の形式でまとめた. 通し番号はゼミ資料と同じものである. ただし, 「性質」はゼミ資料に明示しているか否かによらず, 「定義」「定理」「命題」のどれにも含まれていないかつ大事なものを入れた. また定義は, ゼミ資料中に直接「Def」と書かれていないものでも入っている.
\end{itembox}
\tableofcontents

\clearpage
\section{Review}
    ここは, 高専で習っていない範囲のみを扱う.
    \begin{enumerate}
        \item 行列をブロック分割することで, 見通しがよくなることがある. 例えば, 
        \begin{equation*}
            \begin{pmatrix}
                A & B \\ C & D
            \end{pmatrix}
            \begin{pmatrix}
                F & G \\ H & K
            \end{pmatrix} = 
            \begin{pmatrix}
                AF+BH & AG+BK \\ CF+DH & CG+DK
            \end{pmatrix}
        \end{equation*}
        ただし, $A,B,C,D,F,G,H,K$は行列で, 各々の積が定義できるとする.
        \item $A\bm{e}_i$は$A$の$i$番目の列ベクトル. $\trans\bm{e}_i A$は$A$の$i$番目の行ベクトル.
        \item 行基本変形は, 「正則行列を左から掛けること」である.
        \item 二つの正方行列$A,B$が相似であるとは, ある正則行列$P$を用いて$P^{-1}AP=B$とできること.
        \item 行列式の定義は$|A|=\sum\limits_{\sigma\in S_{n}} \mathrm{sgn}(\sigma) \hspace{1mm} a_{1\sigma(1)}a_{2\sigma(2)}\cdots a_{n\sigma(n)}$
        \item 集合$X,Y$について, $X\subset Y$とは任意の$x\in X$が$x\in Y$を満たすことをいう. $X=Y$とは$X\subset Y$かつ$Y\subset X$である.
        \item 和集合$A\cup B = \{x\mid x\in A または x\in B\}$, 共通部分$A\cap B = \{x\mid x\in A かつ x\in B\}$
        \item $f:X\to Y$が写像とは, 「任意の$X$の元に対してある$Y$の元がただ一つに定まる対応」のこと. 二つの$f,g$ が等しいのは, 任意の$x\in X$に対し$f(x)=g(x)$が成立すること.
        \item $A\subset X$として, 像$f(A)=\{f(x)\mid x\in A\}$.
        \item $B\subset Y$として, 逆像$f^{-1}(B)=\{x\in X\mid f(x)\in B\}$.
        \item $f$が単射とは, 任意の$x,y\in X$に対し, 「$x\neq y\rightarrow f(x)\neq f(y)$」が成り立つこと.
        \item $f$が全射とは, 任意の$y\in Y$に対し, ある$x\in X$があって, $y=f(x)$となること.
    \end{enumerate}
\clearpage
\section{線形空間}
    \begin{definition}
        $V$が$\K$上の\textbf{線形空間}(ベクトル空間)とは
        \begin{enumerate}\renewcommand{\labelenumi}{(\Roman{enumi})}
            \item $\bm{x},\bm{y}\in V$に対し, 和$\bm{x}+\bm{y}\in V$が定まり
            \begin{enumerate}\renewcommand{\labelenumii}{(\arabic{enumii})}
                \item $(\bm{x}+\bm{y})+\bm{z}=\bm{x}+(\bm{y}+\bm{z})$
                \item $\bm{x}+\bm{y}=\bm{y}+\bm{x}$
                \item $\exists\bm{0}\in V,\forall \bm{x}\in V, \bm{x}+\bm{0}=\bm{x}$
                \item $\forall\bm{x}\in V,\exists\bm{x}'\in V, \bm{x}+\bm{x}'=\bm{0}$
            \end{enumerate}
            \item $\bm{x}\in V,a\in\K$に対し, スカラー倍$a\bm{x}\in V$が定まり
            \begin{enumerate}\setcounter{enumii}{4}\renewcommand{\labelenumii}{(\arabic{enumii})}
                \item $(a+b)\bm{x}=a\bm{x}+b\bm{x}$
                \item $a(\bm{x}+\bm{y})=a\bm{x}+a\bm{y}$
                \item $(ab)\bm{x}=a(b\bm{x})$
                \item $1\bm{x}=\bm{x}$
            \end{enumerate}
        \end{enumerate}
    \end{definition}
    \begin{definition}
        $\K$上の線形空間$V,V'$に対し, $T:V\to V'$として
        \begin{equation*}
            T(\bm{x}+\bm{y}), T(a\bm{x})=aT(\bm{x})\quad (a\in\K,\bm{x},\bm{y}\in V)
        \end{equation*}
        を満たすとき, $T$を$V$から$V'$への\textbf{線形写像}という.
    \end{definition}
    \begin{definition}
        $\varphi:V\to V'$が線形写像かつ全単射なら, $V$から$V'$への\textbf{同型写像}という. もし同型写像が存在すれば, $V,V'$は互いに\textbf{同型}といい, $V\simeq V'$とかく.
    \end{definition}
    以下, 特に明記しない限り, $V$を$\K$上の線形空間とする.
    \begin{definition}
        $V$のベクトル$\bm{a}_1,\dots,\bm{a}_n$に対し
        \begin{equation*}
            c_1\bm{a}_1+\cdots+c_n\bm{a}_n\quad (c_i\in\K)
        \end{equation*}
        を$\bm{a}_1,\dots,\bm{a}_n$の\textbf{一次結合}という. また, 
        \begin{equation*}
            c_1\bm{a}_1+\cdots+c_n\bm{a}_n=\bm{0}
        \end{equation*}
        なる$\bm{a}_1,\dots,\bm{a}_n$の関係を\textbf{一次関係}という. 
    \end{definition}
    \begin{property}
        $\bm{a}_1,\dots,\bm{a}_n$の間には少なくとも一つは一次関係が存在する. それは$c_1=\cdots=c_n=0$なるときの関係(\textbf{自明な一次関係}).
    \end{property}
    \begin{definition}
        $\bm{a}_1,\dots,\bm{a}_n$の間に自明でない一次関係が存在するとき\textbf{一次従属}といい, 存在しないとき\textbf{一次独立}という.
    \end{definition}
    \begin{property} 
        一次独立・従属の性質.
        \begin{enumerate}
            \item $\bm{a}_1,\dots,\bm{a}_n$が一次従属であることと, $\bm{a}_1,\dots,\bm{a}_n$のうちある一つのベクトルが他の$n-1$個のベクトルの一次結合でかけることは同値.
            \item $\bm{a}_1,\dots,\bm{a}_n$が一次独立で, $\bm{a}$が$\bm{a}_1,\dots,\bm{a}_n$の一次結合でかけなければ, $\bm{a}_1,\dots,\bm{a}_n,\bm{a}$は一次独立.
            \item $\bm{c}$が$\bm{b}_1,\dots,\bm{b}_l$の一次結合で, 各$\bm{b}_i$が$\bm{a}_1,\dots,\bm{a}_k$の一次結合なら, $\bm{c}$も$\bm{a}_1,\dots,\bm{a}_k$の一次結合.
        \end{enumerate}
    \end{property}
    \clearpage
    \begin{definition}
        $V$に有限個のベクトルがあって, $V$の任意のベクトルがこれら有限個のベクトルの一次結合でかけるとき, $V$を\textbf{有限次元}という. 有限次元でないとき\textbf{無限次元}という. (基本的には有限次元を仮定.)
    \end{definition}
    \begin{definition}
        $V$の有限個のベクトル$\bm{e}_1\dots,\bm{e}_n$が次の二条件を満たすとき, \textbf{基底}という.
        \begin{enumerate}
            \item $\bm{e}_1\dots,\bm{e}_n$は一次独立.
            \item $V$の任意の元は, $\bm{e}_1\dots,\bm{e}_n$の一次結合でかける.
        \end{enumerate}
    \end{definition}
    \begin{property}
        $V$の元を$\bm{e}_1\dots,\bm{e}_n$の一次結合で表したとき, その各$\bm{e}_i$の係数は一意に定まる.
    \end{property}
    \begin{proposition}
        $\K^n$において$n$個より多くのベクトルは一次従属.
    \end{proposition}
    \begin{property}
        $V$が$n$個の基底を持てば, $V\simeq\K^n$
    \end{property}
    \begin{definition}
        $V$の基底が含むベクトルの数を\textbf{次元}といい, $\dim V$とかく. $V=\{\bm{0}\}$のときは$\dim V=0$とする.
    \end{definition}
    \begin{property}
        次元についての諸性質.
        \begin{itemize}
            \item $\dim V$より多いベクトルの組は必ず一次従属.
            \item 次元が等しいベクトル空間は互いに同型.
        \end{itemize}
    \end{property}
    \begin{property}
        $m\times n$行列$A\in M_{nm}(\K)$の解空間$\{\bm{x}\mid A\bm{x}=\bm{0}\}$の次元は$n-\rank A$.
    \end{property}
    \begin{definition}
        ベクトル空間$V$の二つの基底$\E=<\bm{e}_1,\dots,\bm{e}_n>,\F=<\bm{f}_1,\dots,\bm{f}_n>$について, $\E\to\F$の\textbf{基底の取り換え行列}$P$は以下で定義される.
        \begin{equation*}
            (\bm{f}_1,\dots,\bm{f}_n)=(\bm{e}_1,\dots,\bm{e}_n)P
        \end{equation*}
        もしくは次のように定義してもよい. $V$の任意の元を$\E,\F$それぞれの基底で表したときの各係数を縦に並べたもの$\trans(x_1,\dots,x_n),\trans(y_1,\dots,y_n)$について
        \begin{equation*}
            \begin{pmatrix}
                x_1\\\vdots\\x_n
            \end{pmatrix}
            = P
            \begin{pmatrix}
                y_1\\\vdots\\y_n
            \end{pmatrix}
        \end{equation*}
    \end{definition}
    \begin{definition}
        $\K$上の$V$の部分集合$W$が同じ演算で$\K$上の線形空間となるとき, $W$を$V$の\textbf{(線形)部分空間}という.
    \end{definition}
    \begin{property}
        $W$が部分空間である必要十分条件は, 次の三つである.
        \begin{enumerate}
            \item $W\neq \varnothing$.
            \item $\bm{x},\bm{y}\in W\rightarrow \bm{x}+\bm{y}\in W$.
            \item $\bm{x}\in W,c\in \K\rightarrow c\bm{x}\in W$.
        \end{enumerate}
    \end{property}
    \begin{proposition}
        $W_1,W_2$が$V$の部分空間なら, $W_1\cap W_2$も部分空間.
    \end{proposition}
    \begin{definition}
        空でない$V$の部分\underline{集合}$S\subset V$に対し,
        \begin{equation*}
            \Span(S) = \{c_1\bm{x}_1+\cdots+c_{k}\bm{x_k}\mid c_i\in\K,\bm{x}_i\in S\}
        \end{equation*}
        は$V$の部分空間. これを\textbf{$S$によって張られる(生成される)部分空間}という.
    \end{definition}
    \clearpage
    \begin{proposition}
        $W_1,W_2$(部分空間)に対して
        \begin{equation*}
            W_1+W_2=\{\bm{x}_1+\bm{x}_2\mid \bm{x}_1\in W_1,\bm{x}_2\in W_2\}
        \end{equation*}
        を$W_1,W_2$の\textbf{和空間}という. これは部分空間である.
    \end{proposition}
    \begin{definition}
        $f:V\to V'$を線形写像とする. $f$の\textbf{像}を$\Im f=f(V)$とかく. また, $f$の\textbf{核}を$\Ker f = f^{-1}(\{\bm{0}\})$とかく. これらも部分空間である.
    \end{definition}
    \begin{theorem}[次元定理]
        $f:V\to V'$を線形写像とする. このとき次が成り立つ.
        \begin{equation*}
            \dim V=\dim \Im f+\dim\Ker f
        \end{equation*}
    \end{theorem}
    \begin{proposition}
        $W_1,W_2$が$V$の部分空間なら
        \begin{enumerate}
            \item $W_1\subset W_2\rightarrow \dim W_1\leq \dim W_2$
            \item $W_1\subset W_2$かつ$\dim W_1=\dim W_2\rightarrow W_1=W_2$
        \end{enumerate}
    \end{proposition}
    \Slash
    \begin{property}
        $\Im f$の基底は, $V$の基底$\bm{e}_1,\dots,\bm{e}_n$を$f$で飛ばした$f(\bm{e}_1),\dots,f(\bm{e}_n)$の中にある. よってこの中から一次独立なものを取ればそれは基底.
    \end{property}
    \begin{property}
        線型空間$V$の$\dim V$個の線形独立なベクトルは$V$の基底.
    \end{property}
    \clearpage
\section{線形写像と計量線型空間}
    \textcolor{red}{注. 以下, 混同の恐れがない限り, ベクトルを太字で書かずにプレーンテキストで書く.}
    \begin{definition}
        $V,V'$をベクトル空間, $V$の基底を(対応する同型写像と合わせて)$(\E,\varphi)$, $V'$の基底を$(\E',\varphi')$とする. 線形写像$T:V\to V'$に対して, $T_A=\varphi'\circ T\circ \varphi$は$\K^n\to \K^m$の線形写像. 
        したがって, ある$n\times m$の行列$A$でかけるから, その行列$A$を$(\E,\varphi),(\E',\varphi')$に関する$T$の\textbf{表現行列}という. これは, 次式で与えられる.
        \begin{equation*}
           (Te_1,\dots,Te_n)=(e_1',\dots,e_m')A 
        \end{equation*}
        ただし,$\E=<e_1,\dots,e_n>,\E'=<e_1',\dots,e_m'>$. またこれは$V,V'$の$\E,\E'$に対する``成分''に対しては
        \begin{equation*}
            \begin{pmatrix}
               x_1' \\ \vdots \\ x_m' 
            \end{pmatrix}
            =A\begin{pmatrix}
                x_1 \\ \vdots \\ x_n 
            \end{pmatrix}
        \end{equation*}
        とかける.
    \end{definition}
    \begin{property}
        上記の基底とは別に, 基底$\F,\F'$をとり, $\E\to\F$の取り換え行列を$P$, $\E'\to\F'$の取り換え行列を$Q$とする. このとき, $\E',\F'$に関する$T$の表現行列$B$は
        \begin{equation*}
            B=Q^{-1}AP
        \end{equation*}
    \end{property}
    \begin{proposition}
        $V\to V'$の線形写像$T$に対し, $V,V'$の基底を適当に選べば, その表現行列は標準形$\begin{pmatrix}E_{r} & O \\ O & O\end{pmatrix}$ .
    \end{proposition}
    \begin{definition}
        $T:V\to V'$の\textbf{階数}を$\dim(\Im T)$と定義する.
    \end{definition}
    \begin{proposition}
        $T$の階数は任意の表現行列の階数と等しい.
    \end{proposition}
    \begin{proposition}
        $(m,n)$型行列$A$の階数は$A$の一次独立な列ベクトル(行ベクトル)の最大数に等しい.
    \end{proposition}
    \begin{property}
        $(m,n)$型行列$A$の階数に関する同値な条件.
        \begin{enumerate}
            \item $A$を標準形にした時の0でない行の数.
            \item $A$の0でない小行列式の最大次数.
            \item $A$によって定まる線形写像$T_A$の$\dim(\Im T_A)$.\footnote{$\Im (T_A)$を単に$\Im(A)$と書くこともある.}
            \item $A$の一次独立な列ベクトルの最大数.
            \item $A$の一次独立な行ベクトルの最大数.
        \end{enumerate}
    \end{property}
    \begin{property}
        $T$が線形変換なら, $\E' = \E$として表現行列は考える. このとき基底の取り換え$\E\to \F$によって表現行列は$B=P^{-1}AP$となる. つまり\underline{線形変換の表現行列は互いに相似}.
    \end{property}
    \clearpage
    \begin{definition}
        $\K$上の線形空間$V$が次の公理を満たすとき, $V$を\textbf{計量線型空間}という.\\
        $V$の二元$\bm{x},\bm{y}$に\textbf{内積}と呼ばれる$\K$の元$(\bm{x},\bm{y})$が定まり, 次の性質を満たす.
        \begin{enumerate}
            \item $(\bm{x},\bm{y}_1+\bm{y}_2)=(\bm{x},\bm{y}_1)+(\bm{x},\bm{y}_2),\quad (\bm{x}_1+\bm{x}_2,\bm{y})=(\bm{x}_1,\bm{y})+(\bm{x}_2,\bm{y})$
            \item $(c\bm{x},\bm{y})=c(\bm{x},y),\quad (\bm{x},c\bm{y})=\overline{c}(\bm{x},y)$
            \item $(\bm{x},\bm{y})=\overline{(\bm{y},\bm{x})}$
            \item $(\bm{x,x})\geq 0$, 等号は$\bm{x}=\bm{0}$のときに限る.
        \end{enumerate}
        $\K=\R$のとき\textbf{実計量空間}, $\K=\C$のとき\textbf{ユニタリ空間}という. $\sqrt{(\bm{x},\bm{x})}$を$\bm{x}$の\textbf{長さ}, \textbf{ノルム}といい, $\|\bm{x}\|$とかく.
        $(\bm{x},\bm{y})=0$のとき, $\bm{x},\bm{y}$は\textbf{直交する}という.
    \end{definition}
    \begin{property}
        通常の$\K^n$の内積は$(\bm{x},\bm{y})=\trans\bm{x}\overline{\bm{y}}$である. 特に$\K=\R$ならバールはいらない.
    \end{property}
    \begin{proposition} 大事な不等式2つ.
        \begin{enumerate}
            \item $|(x,y)|\leq \|x\|\|y\|$ \hfill(\textbf{Schwarzの不等式})
            \item $\|x+y\|\leq \|x\|+\|y\|$ \hfill(\textbf{三角不等式})
        \end{enumerate}
    \end{proposition}
    \begin{proposition}
        0でない$x_1,\dots,x_k$が互いに直交するなら一次独立.
    \end{proposition}
    \begin{definition}
        計量空間$V$の$e_1,\dots,e_k$が互いに直交し, かつ大きさが1のとき, それらは\textbf{正規直交系}という. とくに, $V$の基底なら\textbf{正規直交基底}という.
    \end{definition}
    \begin{property}[Schmidtの直交化法]
        基底$a_1,\dots,a_n$から正規直交基底$e_1\dots,e_n$を作る手順は以下のとおりである.
        \begin{enumerate}
            \item $\displaystyle e_1=\frac{1}{\|a_1\|} a_1$とおく.
            \item $i\geq 2$なら
            \begin{equation*}
                e_i'=a_i-(a_i,e_1)e_1-\cdots - (a_i,e_{i-1})e_{i-1}
            \end{equation*}
            \item $\displaystyle e_i=\frac{1}{\|e_i'\|} e_i'$とする. $i< n$なら再び(2)に戻る.
        \end{enumerate}
        このとき, $e_{m}$は$a_1,\dots,a_{m}$の一次結合. ($1\leq m\leq n$)
    \end{property}
    \begin{definition}
        $V,V'$がそれぞれ$\K$上の計量空間で, $\varphi:V\to V'$が同型写像かつ
        \begin{equation*}
            (\varphi(x),\varphi(y))=(x,y)
        \end{equation*}
        を満たすとき, $\varphi$を$V$から$V'$への\textbf{計量同型写像}という. このような写像が存在するとき, $V,V'$は\textbf{計量同型}という.
    \end{definition}
    \begin{property}
        $V$の基底$\E$が\underline{正規直交基底}なら$\K^n$への自然な同型写像$\varphi$は計量同型写像.
    \end{property}
    \begin{proposition}
        $V$の正規直交基底$\E,\F$に対し, $\E\to\F$の基底の取り換え行列$P$はユニタリ行列. ($\K=\R$なら直交行列)
    \end{proposition}
    \begin{definition}
        ユニタリ[実計量]空間$V\to V$の計量同型写像を$V$の\textbf{ユニタリ変換}[\textbf{直交変換}]という.
    \end{definition}
    \begin{proposition}
        $V$のユニタリ変換$T$の任意の正規直交基底に関する表現行列はユニタリ行列. 逆にある正規直交基底で$T$の表現行列なら$T$はユニタリ変換.
    \end{proposition}
    \Slash
    \begin{property}
        計量空間$V$の部分空間$W$に対し, $W$のすべての元と直交するような$V$の元全体を$W$の\textbf{直交補空間}といい$W^{\bot}$で表す.
        \begin{enumerate}
            \item $W^{\bot}$は$V$の部分空間.
            \item $V=W\oplus W^{\bot}$, すなわち, $V$は$W$と$W^{\bot}$の直和である. 
        \end{enumerate}
        なお一般に, ベクトル空間$V$の部分空間$W_1,W_2$について, $V$が$W_1,W_2$の和空間$W_1+W_2$で, $V$の任意のベクトルの$W_1,W_2$のベクトルの和として表す仕方が一意的であるとき, \textbf{直和}であるという.
    \end{property}
    \begin{property}
        線形写像と一次独立性に関する性質.
        \begin{enumerate}
            \item 線形写像$f$が単射であることと$\Ker f=\{\bm{0}\}$が同値である.
            \item 任意の線形写像$f$について, $f(v_1),\dots,f(v_k)$が一次独立なら$v_1,\dots,v_k$は一次独立である.
            \item 単射な線形写像$f$については, $v_1,\dots,v_k$が一次独立なら$f(v_1),\dots,f(v_k)$は一次独立である.
        \end{enumerate}
    \end{property}
    \clearpage
\section{固有値(計量無し)}
    \begin{definition}
        $V$を$\K$上のベクトル空間, $T$を$V$の線形変換とする.
        \begin{equation*}
            T\bm{x}=\alpha \bm{x}\quad (\alpha\in\K)
        \end{equation*}
        なる$\bm{x}\neq\bm{0}$を$T$の\textbf{固有ベクトル}といい, $\alpha$を$T$の\textbf{固有値}という. $\alpha$に対応する固有ベクトル全てに$\bm{0}$を入れた集合$W_\alpha$を$T$の$\alpha$に対する\textbf{固有空間}という. 固有空間は$T$不変である.
    \end{definition}
    \begin{definition}
        $A$が$n$次正方行列のとき, $A$によって決まる\underline{$\C^n$}の線形変換$T_A$の, 固有値, 固有ベクトル, 固有空間をそれぞれ$A$の固有値, 固有ベクトル, 固有空間という.
    \end{definition}
    \begin{property}
        $T$が$\C$上のベクトル空間の線形変換なら, $T$の表現行列$A$の固有値は$T$の固有値.
    \end{property}
    \begin{proposition}
        $T$の相異なる固有値に対する固有ベクトルは一次独立.
    \end{proposition}
    \begin{proposition}
        $\K$上のベクトル空間$V$の線形変換が, 適当な基底に関して対角行列で表現されるためには
        \begin{equation*}
            V=W_1\oplus\cdots\oplus W_k
        \end{equation*}
        が必要かつ十分である. つまり「固有ベクトルからなる$V$の基底が作れること」が必要かつ十分.
    \end{proposition}
    \begin{proposition}
        $n$次正方行列$A$が適当な正則行列$P$によって\textbf{対角化可能}である必要十分条件は$n$個の一次独立な固有列ベクトルが存在することである. 特に, 一次独立な固有ベクトルを$\bm{p}_1,\cdots,\bm{p}_n$として$P=(\bm{p}_1\cdots\bm{p}_n)$とすると
        \begin{equation*}
            P^{-1}AP=\begin{pmatrix}
                \alpha_1 & 0 & \cdots & 0 \\
                0 & \alpha_2 & \cdots & 0 \\
                \vdots & \vdots & \ddots & \vdots \\
                0 & 0 & \cdots & \alpha_n
            \end{pmatrix}
        \end{equation*}
    \end{proposition}
    \begin{definition}
        $A$を$n$次行列とする. $\Phi_A(x)=\det(xE-A)$を$A$の\textbf{特性多項式}, \textbf{固有多項式}という. $\Phi_A(x)=0$を$A$の\textbf{特性方程式}, \textbf{固有方程式}という. 特性方程式の根を$A$の\textbf{特性根}という.
        $T$の\underline{任意の}表現行列に関する特性多項式, $\cdots$を$T$の\textbf{特性多項式}, $\cdots$という.
    \end{definition}
    \begin{property}
        基底の取り方で$T$の特性多項式は変わらない.
    \end{property}
    \begin{property}[代数学の基本定理]
        複素係数$n$次方程式は重複度込めてちょうど$n$個の解をもつ.
    \end{property}
    \begin{proposition}
        行列$A$(or $\C$上の線形変換$T$)の固有値は, $A$(or $T$)の特性根と一致する.
    \end{proposition}
    \begin{proposition}[対角化可能である必要十分条件]
        行列$A$が対角化可能である必要十分条件は, $A$の各特性根$\alpha$に対する固有空間の次元が$\alpha$の重複度に一致することである. 特に, 特性根がすべて相異なるなら$A$は対角化可能である.
    \end{proposition}
    \clearpage
\section{固有値(計量有り)}
    \begin{definition}
        実計量空間$V$の線形変換$T$が, 任意の$\bm{x},\bm{y}\in V$に対し
        \begin{equation*}
            (T\bm{x},\bm{y})=(\bm{x},T\bm{y})
        \end{equation*}
        なら, $T$を\textbf{対称変換}という.
    \end{definition}
    \begin{proposition}
        $T$が対称変換なら, $V$の任意の正規直交基底に関する$T$の表現行列は, 実対称行列. 逆にある正規直交基底に関する$T$の表現行列が実対称行列なら, $T$は対称変換.
    \end{proposition}
    \begin{proposition}
        対称変換の特性根はすべて実数.
    \end{proposition}
    \begin{proposition}
        実計量空間$V$の対称変換$T$の相異なる固有値を$\beta_1,\dots,\beta_k$, 対応する固有空間を$W_1,\dots,W_k$とすると, $W_1,\dots,W_k$は互いに直交し
        \begin{equation*}
            V=W_1\oplus\cdots\oplus W_k
        \end{equation*}
    \end{proposition}
    \begin{proposition}[系]
        実計量空間$V$の線形変換$T$が適当な正規直交基底に関したい対角行列で表現される必要十分条件は, $T$が対称変換であることである.
    \end{proposition}
    \begin{theorem}
        実正方行列$A$に対し$P^{-1}AP$が対角行列であるような直交行列$P$が存在するためには, $A$が対称行列であることが必要かつ十分である.
    \end{theorem}
    \begin{definition}
        $n$個の変数$x_1,\dots,x_n$に対する実係数斉二次式を\textbf{二次形式}という. 一般式は次で与えられる.
        \begin{equation*}
            F(x_1,\cdots,x_n)=\sum_{i,j=1}^{n}a_{ij}x_ix_j,\quad (a_{ij}=a_{ji})
        \end{equation*}
        対称行列と列ベクトルを用いれば次のように書ける.
        \begin{equation*}
            F(\bm{x})=\trans\bm{x}A\bm{x}
        \end{equation*}
    \end{definition}
    \begin{definition}
        二次形式は適当な直交行列$P$による変数変換$\bm{x}=P\bm{y}$で次の\textbf{標準形}\footnote{ゼミ資料では係数まで丸め込んで標準形といったが, 固有値は陽に出しておくのが一般的なのでこちらを採用する.}に変形できる.
        \begin{equation*}
            F(\bm{x})=G(\bm{y})=\alpha_1y_1^2+\cdots+\alpha_n y_n^2
        \end{equation*}
        ただし, $\alpha_1,\dots,\alpha_n$は$A$の固有値.
    \end{definition}
    \begin{theorem}[Sylvesterの慣性法則]
        任意の正則変換によって二次形式を標準形に変形したとき、正の項の数、負の項の数は変換によらず一意に決まる。
    \end{theorem}
    \begin{definition}
        $\bm{0}$でない任意の$\bm{x}$に対し, $f(\bm{x})>0$が成立するとき, 二次形式$f(\bm{x})$は\textbf{正値}であるという.
    \end{definition}
    \begin{proposition}
        $F(\bm{x})$が正値であることと$p=n$ ($A$の固有値がすべて正)は同値.
    \end{proposition}
    \begin{proposition}
        $F(\bm{x})$が$\|\bm{x}\|=1$という条件のものとでとる最大値は, 二次形式を定める行列の固有値の最大値と等しい.
    \end{proposition}
    \begin{proposition}
        二次形式$F(\bm{x})=\trans\bm{x}A\bm{x}$が正値ならば$|A_k|>0\hspace{1mm}(k=1,2,\dots,n)$. 逆も成り立つ. ここで
        \begin{equation*}
            A_k=\begin{pmatrix}
                a_{11} & \cdots & a_{1k} \\
                \vdots & \ddots & \vdots\\
                a_{k1} & \cdots & a_{kk}
            \end{pmatrix}
        \end{equation*}
        である.
    \end{proposition}
    \begin{theorem}[三角化定理]
        任意の行列は\textbf{三角化可能}, つまり, ある正則行列$P$をとって
        \begin{equation*}
            P^{-1}AP = \begin{pmatrix}
                \alpha_1 & \multicolumn{3}{c}{\smash{\raisebox{-2.5ex}{\Huge *}}} \\
                0 & \alpha_2 & & \\
                \vdots & \vdots & \ddots & \\
                0 & 0 & \cdots & \alpha_n
            \end{pmatrix}
        \end{equation*}
        とできる. ここで$\alpha_1,\dots,\alpha_n$は$A$の固有値である.
    \end{theorem}
    \begin{proposition}
        任意の正方行列$A$とその固有値を$\alpha_1,\dots,\alpha_n$として, 次が成り立つ.
        \begin{enumerate}
            \item $|A|=\alpha_1\alpha_2\cdots\alpha_n$
            \item $\tr A=\alpha_1+\alpha_2+\cdots+\alpha_n$
        \end{enumerate}
    \end{proposition}
    \begin{definition}
        行列$A$の\textbf{多項式}とは$E,A,A^2,\cdots,A^k$の一次結合のこと. そして, `普通の'多項式
        \begin{equation*}
            f(x)=a_{k}x^k+a_{k-1}x^{k-1}+\cdots+a_{1}x+a_{0}
        \end{equation*}
        に対して
        \begin{equation*}
            f(A)=a_{k}A^k+a_{k-1}A^{k-1}+\cdots+a_{1}A+a_{0}E
        \end{equation*}
        を多項式$f(x)$の$A$での\textbf{値}という.
    \end{definition}
    \begin{theorem}[Cayley-Hamiltonの定理]
        任意の行列$A$はその特性多項式の零点. すなわち
        \begin{equation*}
            \Phi_A(A)=O
        \end{equation*}
    \end{theorem}
    \Slash
    \begin{property}[Frobeniusの定理]
        $g(t)$を任意の多項式, $\alpha_1,\dots,\alpha_n$を$n$次正方行列$A$の固有値とする. このとき, $g(A)$の固有値は$g(\alpha_1),\dots,g(\alpha_n)$である.
    \end{property}
    \begin{property}[次元公式]
        $W_1,W_2$をベクトル空間$V$の部分空間とする. この時次式が成り立つ.
        \begin{equation*}
            \dim(W_1+W_2)=\dim W_1+\dim W_2 - \dim(W_1\cap W_2)
        \end{equation*}
    \end{property}
    \clearpage
    \section{Tips}
        ここでは覚えておいたほうがいいことをまとめてあげる. 重複も問わない.
        \subsection{行列の相似}
            相似な行列で不変なものはつぎである.
            \begin{itemize}
                \item 階数
                \item 行列式
                \item トレース
                \item 固有値
                \item 固有多項式
            \end{itemize}
        \subsection{連立方程式の理論}
            $m\times n$の係数行列$A$と$n$項列ベクトル$\bm{x}$について, 斉次方程式$A\bm{x}=\bm{0}$は$n-\rank A$個の自明でない解をもつ. つまり, 解空間$\{\bm{x}\mid A\bm{x}=\bm{0}\}$の次元は$n-\rank A$.

            一般の連立方程式$A\bm{x}=\bm{b}$に対して, 拡大係数行列は$\tilde{A}=[A \mid \bm{b}]$と定義されるが, このときこの連立方程式が解をもつ必要十分条件は$\rank A=\rank \tilde{A}$である.
        
        \subsection{トレース(trace)}
            $\displaystyle\tr A=\sum_{i=1}^{n}a_{ii}$を$A$のトレース(対角和, 跡)という. $A$の対角成分の和である. 性質としては次の三つを抑えておけばよい.
            \begin{equation*}
                (1)\tr(A+B)=\tr (A)+\tr (B),\quad (2)\tr(kA)=k\hspace{1mm}\tr A\quad (3)\tr(AB)=\tr(BA).
            \end{equation*}
        \subsection{ユニタリ行列に関する性質}
            $A^*\equiv \trans\overline{A}$とする. $AA^*=A^*A=E$なる$A$をユニタリ行列という. 特に$A$が実行列なら直交行列という. ユニタリ行列について, 次の四つはすべて同値である. ただし, $(\bm{x},\bm{y})=\trans\bm{x}\overline{\bm{y}}$ (標準内積)とする.
            \begin{equation*}
                (1)Aはユニタリ行列,\quad(2)(A\bm{x},A\bm{y})=(\bm{x},\bm{y}),\quad(3)\|A\bm{x}\|=\|\bm{x}\|,\quad(4)Aの各列ベクトルは正規直交基底
            \end{equation*}
        \subsection{三角行列について補足}
            上下によらず三角行列の行列式は対角成分の積である. また, 三角化すればわかるように$A$の固有値が$\alpha_1\dots,\alpha_n$なら$A^{k}$の固有値は$\alpha_1^k,\cdots,\alpha_n^k$である. (Frobeniusの定理からもわかる.)
\end{document}